\documentclass[11pt]{article}
\usepackage[margin=1in]{geometry}
\usepackage[T1]{fontenc}
\usepackage{lmodern,amsmath,amssymb,booktabs,graphicx}
\usepackage[colorlinks=true,linkcolor=blue,urlcolor=blue]{hyperref}
\setlength{\parindent}{0pt}\setlength{\parskip}{5pt}
\setlength{\emergencystretch}{2em}
\title{RoCE-K: target protection, model approximation and design stability}
\author{}\date{4 October 2026}
\begin{document}\maketitle

\textbf{Summary.} The target-protection change gives a useful reduction beyond
range estimation alone, while improving the target-only comparator in parallel.
A separate stress test exposes a real failure of the uncorrected conditional
source certificate under nonlinear invalid means. Correcting variance centering
and its concentration bound restores the guarantee within a stated model class.
The design study makes the large-$K$ fallback problem concrete. These are
separate conclusions; favorable population coverage is not used to conceal a
failed conditional source interval.

\section{Accounting and unchanged scientific targets}

There are 7,200 outcome setting-specific repetitions: 12 target-protection
settings, eight random-design approximation settings and 16 fixed-design stress
settings, each with 200 repetitions and 1,000 patients per site. Methods share
data within each repetition, and approximation amplitudes share seeds within
source-count settings. Thus these are not 7,200 mutually independent datasets.
A further 7,800 design-only repetitions vary source sample sizes. Paired broader-
model reanalyses of 4,800 existing repetitions are not counted again.

The random-design DGP retains the preceding bounded continuous-feature model,
source covariate tilt .4, four outcome signal slopes, and population TATE .10.
Target dimensions are $p=4,10$, source counts are $64,256$, and target CATE
heterogeneity is $h=0,.15$. Constant-effect checks also use $K=2048$; all-valid
controls use $K=256$. Each arm is guaranteed at least $3K/4$ valid sources.

The approximation panel leaves the means at target and valid sources linear and adds
$\eta x_1x_2$ to treated means only at invalid sources, retaining their weak
intercept shift. The final grid is $\eta=0,.05,.10,.15$. A check of the probability support
rejected the original $\eta=.2$ setting because some support corners
have negative probabilities. That aborted attempt remains archived; there is
no probability clipping. Unchanged completed cells were reused only after
setting equality and exact replay checks.

All feasible procedures use observed summaries and declared model assumptions.
True coefficients, variances and target ranges are used only for audit and
coverage evaluation. Both conditional and population intervals, source bands,
noise radii, dispersion radii and composition radii are saved.

\section{Target protection with a fair target-only comparator}

A single conditional sub-Gaussian confidence ellipsoid for target treatment-
effect slopes supplies an upper CATE range and an upper empirical CATE variance.
The range-only method plugs the former into Hoeffding; the empirical Bernstein
method uses both. A simultaneous-minimum variant spends half its composition
error on each inequality. Every new source-assisted variant uses identical
conditional source and target bands, verified directly from saved arrays.

The allocation is .005 for coefficient uncertainty, .01 for composition,
.01 each for source and target outcome noise, and .015 for dispersion.
Target-only uses the same coefficient event and .0225 each for its other two
components. The table's universal control uses this matched allocation; v19's
allocation is retained separately in the numerical files. The coverage column
refers to the source-assisted Bernstein interval.

\begin{center}\small
\resizebox{\linewidth}{!}{\input{target_table}}
\end{center}

Range-only protection helps at $p=4$ but is typically capped at 2 at $p=10$.
Variance-sensitive protection improves both dimensions. At $p=4,K=256,h=.15$,
mean length falls from .2615 to .1584, versus .2853 for the simultaneously
improved protected target-only method. Ordinary target-normal remains shorter
at .1214 in that setting. At $p=4,K=2048,h=0$, the new .1119 interval is shorter
than the ordinary target-normal .1228 interval, without supplying the true
constant-effect structure to the algorithm. This improvement is not universal:
at $p=10,K=2048$ the new .1443 interval still exceeds the normal .1234 reference.

\begin{figure}[ht]\centering
\includegraphics[width=\linewidth]{target_protection.pdf}
\caption{Constant-effect weak-bias settings, 200 repetitions each. Both the
source-assisted and protected target-only methods use the new observed-data
certificate. Ordinary normal intervals remain a separate reference.}
\end{figure}

The main design proves that, for fixed dimension and well-conditioned target
designs, the Bernstein composition allowance has order
$\sigma_\tau n_t^{-1/2}+n_t^{-1}$. At constant CATE it is second order,
$O_p(n_t^{-1})$, without telling the algorithm that the effect is constant.
This is a model-specific rate, not a general high-dimensional optimality claim.

\section{Why invalid-source nonlinearity requires a real correction}

For $m_j=E(W_j\mid\mathcal D)$ and $u_j=(I-H_j)m_j$, the leave-out variance
bias is $\zeta_j=\sum_i d_{ji}m_{ji}u_{ji}$. Consequently
\[
 E(\widehat D\mid\mathcal D)=D-a_K\sum_j \zeta_j.
\]
Positive variance bias can suppress the dispersion certificate. The repair
therefore changes both its centering allowance and the MGF variance proxy.
If valid sources are themselves approximate, their mean-bias allowance must
also be included; adding only one of these terms is insufficient.

The random-design panel compares the uncorrected method with a fixed treated-
arm approximation envelope .2 and with the universal bound implied by bounded
outcomes. In the latter, $\|u_{ja}\|_2\le\sqrt{n_{ja}}/2$, and at most
$K-g_a$ sources per arm need this bound. It does not require a supplied
nonlinearity amplitude, but still assumes a linear target and guaranteed valid
sources. At $K=256,\eta=.15$, population lengths are .1592 (uncorrected), .1629
(declared envelope), .1739 (bounded invalid means), and .2848 (improved protected
target-only). The observed conditional coverage in those source rows is 1;
that alone does not validate the uncorrected nonlinear procedure.

\subsection{A fixed-design counterexample isolates the issue}

Each treatment arm has 500 patients at $x=-1,0,1$ in proportions $1/4,1/2,1/4$.
The working basis is $(1,x)$. Valid treated means equal .5 and controls .4;
one quarter of sources have treated means
\[
 .5+\delta_K+\eta(x^2-.5),\qquad
 \delta_K=2\sqrt{.5/1000}\,K^{-1/4}.
\]
The grouped-binomial simulator exactly reproduces the patient-level sufficient
statistics; it is not a Gaussian approximation or a fit to billions of stored
patient records. Literal patient calculations were independently checked.
This is a conditional-design experiment, separate from the i.i.d. random-design
population panel.

At $K=262144$:
\begin{center}\small
\resizebox{\linewidth}{!}{\input{projection_table}}
\end{center}

At $\eta=.4$, the original conditional coverage is zero while population
coverage is 1 after its wide range-2 enlargement. Conditional width is only
.000402 without correction. The declared-envelope repair has coverage 1 and
width .005498. The broader bounded-invalid version has combined conditional
coverage .990 and width .025698; its source band itself covers in every
repetition. Its two combined-band misses occur when the target band removes
truth from the intersection, an event included in the confidence allocation.
A wider source band need not yield nested reported intersections.

\begin{figure}[ht]\centering
\includegraphics[width=\linewidth]{projection_stress.pdf}
\caption{Fixed-design nonlinear-invalid stress at $\eta=.4$. Conditional
coverage and width are shown separately from population enlargement. The two
repairs use narrower declared and broader bounded-invalid model classes.}
\end{figure}

The manuscript provides the corresponding asymptotic counterexample: true
weak-shift dispersion vanishes, but the uncorrected variance-centering bias
stays positive at fixed per-site design. Its nominal upper bound can collapse
to zero. This invalidates that conditional extension, rather than proving an
unconditional lower bound for every nonlinear method. The universal repair
also has a cost: its current source radius can stop improving with $K$ at an
$O(n_s^{-1/2})$ scale. The exact-linear fourth-root rate is not automatically
inherited by this broader class.

\section{Random-design growth and fallback}

The design-only panel varies $n_s=120,200,1000$, $p=10,45$, $K=4,16,64$ and
logistic treatment-allocation strength $0,2.5$, keeping $n_t=1000$. Three extra
$p=4$ settings increase $(n_s,K)$ together. Every design failure is retained.

\begin{center}\small
\resizebox{\linewidth}{!}{\input{design_table}}
\end{center}

At $n_s=120,p=45$ and strength 2.5, any-source failure rises from .45 at $K=4$
to 1 at $K=64$. Raising $n_s$ to 200 removes observed failures in that cell.
The count-adjusted eligibility column asks only whether discarding failed
sites would leave $g_a'=\max(0,g_a-m)>0$; no alternative CI was fitted there.
Eligibility is not evidence of a precision gain.

The main text derives the sufficient bound
\[
 2(K+1)d\exp\{-n_{\min}\kappa/(8L_d)\}
\]
for the lower Gram-eigenvalue event, plus explicit leverage and weight bounds.
A population upper-eigenvalue assumption and a second Chernoff term control
the numerical condition guard when dimension grows. It can be
vacuous in finite samples and is reported as such. The table's design diagnostic
uses the source-only version of this union bound. The growing-$n_s$ trajectory
shows a regime where it becomes informative, although 200 trials cannot
empirically verify extremely small probabilities.

Unconditional coverage follows from the protected fallback even when failures
are common. Useful width and expected width need additional conditions. In
particular, an $O_p$ width bound is insufficient: the manuscript now controls
$E\sqrt{U_D}$ and includes the clipped-interval contribution
$2P(\text{bad design})$. That probability must be controlled on the desired
width scale, not merely shown to tend to zero.

\section{Verification and remaining work}

The research suite passes 56 tests, including exact non-Gaussian coefficient
and dispersion checks, direct leave-out calculations, agreement between grouped and patient
calculations, probability-support checks and scalar Chernoff checks. Final source
replays preserve the v19 reference and representative new intervals. All old
methods in the 24 broader-model paired cells match their stored intervals
exactly. No failure of the target coefficient certificate was observed in the main
random-design panels.

The next task is to improve the usefulness of the broad model-error certificate
and test an implemented design-only source-removal policy. A nonlinear or
sparse target model still needs its own valid fitting-remainder and composition
certificates. Current-paper RoCE/ENAR defaults and all previous results remain
unchanged.

\end{document}
